\documentclass[11pt,a4paper]{article}
\usepackage[utf8]{inputenc}
\usepackage{amsmath}
\usepackage{amsfonts}
\usepackage{hyperref}
\usepackage{geometry}
\geometry{margin=1in}

\title{Technical Note: Controller Prior Bias Correction ($\Phi_{default}$)}
\author{FSA-v5 Research Engineering}
\date{May 10, 2026}

\begin{document}

\maketitle

\section{Introduction}
In the FSA-v5 control framework, the bimodal stimulus plan $\Phi(t) = (\Phi_B, \Phi_S)$ is optimized using a Gaussian Radial Basis Function (RBF) basis. The optimization is performed over a vector of coefficients $\theta \in \mathbb{R}^{2 \cdot n_{anchors}}$. This note documents the correction of the baseline stimulus intensity $\Phi_{default}$ from $1.0$ to $0.3$ and explains the physiological and mathematical necessity of this change.

\section{Mathematical Implementation}
The stimulus at any time $t$ is decoded from the optimized coefficients $\theta$ using a logistic transformation centered on a bias constant $c_\Phi$. 

\subsection{Bias Precomputation}
The constant $\Phi_{default}$ is converted into the logit-space bias $c_\Phi$ in the \texttt{FSAv5ControlGPUTarget} constructor:
\begin{equation}
    c_\Phi = \ln \left( \frac{\Phi_{default} / \Phi_{max}}{1 - \Phi_{default} / \Phi_{max}} \right)
\end{equation}
where $\Phi_{max} = 3.0$ is the saturation limit.

\subsection{Stimulus Decoding}
Inside the GPU kernel, the stimulus for the aerobic channel $\Phi_B$ (and similarly for strength $\Phi_S$) is calculated as:
\begin{equation}
    \Phi_B(t) = \frac{\Phi_{max}}{1 + \exp\left( -\left[ c_\Phi + \sum_{j=1}^{n_{anchors}} \theta_j \cdot R_j(t) \right] \right)}
\end{equation}
where $R_j(t)$ are the RBF basis functions. 

\section{Rationale for Correction}

\subsection{Physiological Baseline vs. Overtraining}
In the FSA-v5 non-linear dynamics:
\begin{itemize}
    \item \textbf{$\Phi \approx 0.3$:} Represents a physiologically sustainable baseline training load. At this level, the system typically maintains Alertness ($A$) within the healthy basin of attraction.
    \item \textbf{$\Phi = 1.0$:} Represents a heavy overtraining load. Sustained training at this intensity leads to rapid accumulation of fatigue ($F$) and high risk of autonomic collapse (transitioning to the $A=0$ basin).
\end{itemize}

\subsection{Optimization Search Efficiency}
The Sequential Monte Carlo (SMC$^2$) optimizer initializes its search with $\theta \approx 0$. 
\begin{itemize}
    \item \textbf{Old Behavior ($\Phi_{default} = 1.0$):} The search was centered on the "Overtraining" regime. Particles had to perform significant work to "un-train" the initial guess just to reach a sustainable state.
    \item \textbf{New Behavior ($\Phi_{default} = 0.3$):} The search is now centered on a healthy baseline. The optimizer can immediately explore nuanced deviations (tapering or peaking) from a stable starting point.
\end{itemize}

\section{Conclusion}
By setting $\Phi_{default} = 0.3$, we align the controller's prior bias with the model's actual steady-state baseline. This reduces the computational effort required for the SMC particles to find viable plans and prevents the optimizer from being "trapped" in high-fatigue regions of the state space during the initial tempering levels.

\section*{Code References}
\begin{itemize}
    \item \texttt{bench\_args.jl}: CLI flag \texttt{--ctrl-phi-default} (Default: 0.3).
    \item \texttt{gpu\_control\_v5.jl}: Implementation of $c_\Phi$ and the sigmoid decoder.
    \item \texttt{bench\_controller.jl}: Propagation of the default value to the kernel.
\end{itemize}

\end{document}
